\section{Conclusion}
\label{sec:conclusion}

We have introduced StabilitySection, an analysis framework for measuring the
cross-census stability of GeoSection-optimal redistricting maps, and the
Census Stability Score (CSS), a composite metric that quantifies this stability
along three dimensions: partisan seat outcomes, natural ratio selection, and
proportionality gap magnitude.

The central claim of this paper is that census stability is a qualitatively
distinct and legally stronger form of algorithmic robustness than within-year
seed stability.
B.7 establishes that GeoSection finds the same answer regardless of random
initialisation.
B.15 establishes that it finds the same answer regardless of which decennial
census snapshot is used as input.
The combination is the strongest available form of geographic determinism:
the natural partition of a state is determined by the physical geography---
not by random chance, not by partisan composition, not by decennial population
dynamics.

\paragraph{Key findings.}
From the complete 2020 GeoSection baseline (44 multi-district states) and the
available 2000 sweep data (47 states), we establish:

\begin{itemize}
\item \textbf{Twelve states expect CSS $\geq 0.90$} (full three-census):
  the Midwest row (IA, KS, NE, IN, MO), the traditional South (AL, AR, LA, KY,
  TN), and New England (ME, WV, RI).
  These states' geographic structures have been stable for two decades; their
  natural ratios are anchored by physical features---agricultural uniformity,
  mountain topography, fixed city boundaries---that do not change with
  decennial population counts.

\item \textbf{Seventeen states expect CSS $0.70$--$0.90$} (moderate stability):
  the industrial Midwest (IL, PA, MI, OH), the growing South (VA, NC, SC),
  and the Pacific Coast (OR, WA).
  These states have stable first-level ratios but may see slight seat-count
  shifts as suburban growth continues.

\item \textbf{Six states expect CSS $< 0.70$} (structurally volatile):
  the Sun Belt growth corridor (TX, AZ, FL, GA, NV, CO).
  These states have undergone the most dramatic population redistribution
  between 2000 and 2020; their metropolitan geographies have reorganised in
  ways that shift the natural ratio.
  This is not an algorithm failure---it is a signal that the geographic
  structure of these states has genuinely changed.

\item \textbf{The Lorenz drift proxy is actionable now}:
  states where $p^*$ drifts by $>0.03$ between available estimates are
  flagged for priority re-evaluation at each redistricting cycle.
  This provides a pre-census early warning that does not require running
  the full GeoSection sweep.
\end{itemize}

\paragraph{CSS as a litigation tool.}
A court presented with a GeoSection-optimal map can be told: ``This is the
natural geographic partition of this state.
We know it is natural because three independent censuses---across twenty years
of population change, two political realignments, and a national demographic
transformation---all converge to the same answer.
Any deviation from this map requires an explanation that is not geography.''

For states with CSS $\geq 0.90$, this is the \emph{strongest available}
procedural-fairness claim in post-\emph{Rucho} redistricting litigation.
It does not require proof of partisan intent (constitutionally difficult).
It does not require a partisan-symmetry measurement (empirically contested).
It requires only the CSS score and the three-census GeoSection run---both
computed from public data using a published algorithm.

\paragraph{Future work.}
\begin{itemize}
\item \textbf{Complete three-census CSS.}
  The 2010 sweep is in progress; completing it will produce the full
  three-census CSS table for all 50 states.
  This is the next session's primary task.

\item \textbf{Jaccard district-assignment stability.}
  Once the 2000-to-2020 tract alignment is implemented, the full district
  Jaccard computation will provide RQ3 (district-assignment stability)
  and RQ5 (Lorenz curve drift analysis) from the plan.

\item \textbf{Population-only perturbation decomposition.}
  Fixing the 2020 graph and substituting 2000 population weights, for the
  10-state decomposition subset, will isolate the contribution of tract
  boundary redesign to instability.

\item \textbf{Predictive model for CSS.}
  Fitting the regression CSS $\sim f(\text{pop growth, boundary growth,
  seat change, } p^*)$ on the complete three-census dataset will enable
  pre-census CSS prediction for 2030 redistricting.

\item \textbf{Temporal sensitivity using ACS 5-year estimates.}
  Testing whether $\Delta p^*$ between consecutive ACS releases predicts
  ratio shifts at the next decennial would allow annual CSS monitoring
  rather than decennial-only evaluation.
\end{itemize}

\paragraph{Code and data availability.}
StabilitySection is implemented in the \texttt{redist} Rust binary.
Cross-census GeoSection runs use \texttt{redist state --year \{2000,2010,2020\}
--partition-mode geosection}.
Sweep outputs are available in the replication archive
(\texttt{research/A.4+replication-materials}).
The CSS computation script is in \texttt{scripts/research/stability\_section/}.
